\section{Требования к тестированию}

\subsection{Тесты бэкенда}

Используются \textbf{pytest} и \textbf{testcontainers} для поднятия реальной PostgreSQL и при необходимости S3-совместимого хранилища (например, MinIO) в тестовом окружении. Покрываются API-эндпоинты, бизнес-логика проверки документов и взаимодействие с БД и хранилищем.

\subsection{Тесты фронтенда}

Используются \textbf{Jest} и \textbf{React Testing Library} (или эквивалентный стек) для модульного и компонентного тестирования интерфейса.

\subsection{Тесты правил проверки документов}

Используются эталонные материалы; тесты должны подтверждать, что корректно оформленные документы проходят соответствующие проверки. Дополнительно проверяются граничные и негативные случаи: повреждённый PDF, неподходящий тип файла, намеренно некорректная структура документа.

\subsection{Тесты безопасности загружаемых PDF}

Отдельно проверяется, что файлы с признаками вредоносного содержимого или типичных PDF-эксплойтов (встроенные исполняемые объекты, JavaScript, опасные действия) отклоняются системой и не сохраняются в хранилище. При необходимости используются специально подготовленные тестовые PDF с известными паттернами риска.
